\chapter*{CHƯƠNG 3. CÔNG NGHỆ VÀ KIẾN TRÚC HỆ THỐNG}
\addcontentsline{toc}{chapter}{CHƯƠNG 3. CÔNG NGHỆ VÀ KIẾN TRÚC HỆ THỐNG}

\section*{3.1 Kiến trúc hệ thống}
Hệ thống quản lý phân bổ phòng ký túc xá được thiết kế theo kiến trúc phân lớp với backend chịu trách nhiệm xử lý nghiệp vụ và truy xuất dữ liệu, frontend thực hiện hiển thị động theo mẫu server-side rendering, cùng các tác vụ ngoại vi đảm nhiệm sinh dữ liệu, kiểm thử và thao tác hàng loạt. Việc lựa chọn kiến trúc này xuất phát từ yêu cầu về tính nhất quán dữ liệu khi thực hiện các toán tử phân bổ theo chu kỳ (allocation cycle), tính mở rộng cho chính sách kế thừa (policy inheritance), và khả năng tự động hóa các kịch bản xử lý vi phạm (violation workflow). Kiến trúc ưu tiên phân tách rõ ràng giữa tầng truy vấn dữ liệu (MongoDB qua Mongoose), tầng API REST (Node.js + Express), và tầng trình bày (EJS), nhằm giảm độ phức tạp đồng thời tạo điều kiện cho nhóm tác vụ sinh dữ liệu mô phỏng và nhóm tác vụ phân phòng tự động vận hành độc lập nhưng vẫn tương tác an toàn với API lõi.

Mô hình luồng xử lý chính bao gồm: client gửi yêu cầu qua trình duyệt, server xác thực và cấp token/session, service thực hiện tính toán điểm ưu tiên (\texttt{priority\_score} allocation) và ghi nhận vào cơ sở dữ liệu, sau đó hệ thống cập nhật trạng thái phân bổ và phản hồi client. Thiết kế này cho phép tách biệt các quy trình dài (batch seeding, migration) khỏi các yêu cầu thời gian thực, đồng thời vẫn bảo đảm nhất quán bằng cách sử dụng các transaction logic ở mức ứng dụng và kiểm tra bất đồng bộ ở các script tác nghiệp.

Để giảm thiểu xung đột khi nhiều tiến trình cùng thao tác trên cùng một tập dữ liệu (ví dụ khi chạy tác vụ phân phòng tự động đồng thời với migration), hệ thống áp dụng các biện pháp đồng bộ hóa ở tầng ứng dụng: sử dụng cơ chế optimistic locking cho các document quan trọng, kiểm tra phiên bản trước khi ghi và thiết kế các job hàng loạt chạy theo hàng đợi. Cách tiếp cận này hạn chế nhu cầu lock cơ sở dữ liệu nặng, giúp duy trì thông lượng cho các truy vấn đọc/ghi trong các đợt phân bổ lớn.

Một điểm then chốt của kiến trúc là khả năng mô hình hóa đầy đủ chu trình nghiệp vụ theo năm học, thay vì chỉ xử lý các thao tác CRUD độc lập. Cụ thể, chức năng tạo cửa sổ học kỳ không chỉ tạo mốc thời gian tiếp nhận hồ sơ, mà còn đóng vai trò mốc điều phối cho toàn bộ pipeline phân bổ: khi học kỳ mới được kích hoạt, hệ thống đồng thời kiểm tra trạng thái policy inheritance, hiệu lực tiêu chí ưu tiên, và khả năng tiếp nhận của từng phòng theo khu vực. Cách tổ chức này giúp tránh tình trạng mỗi module tự diễn giải thời gian theo cách khác nhau, vốn dễ dẫn đến sai lệch kết quả room assignment khi hệ thống vận hành trong thời gian dài.

Bên cạnh đó, kiến trúc còn cần xử lý các tình huống điều chỉnh chính sách sau khi đã có dữ liệu phát sinh. Trên thực tế, khi ban quản lý thay đổi cách tính ưu tiên hoặc điều kiện xét vi phạm, các hồ sơ đã nộp có thể chịu ảnh hưởng dây chuyền. Vì vậy, hệ thống sử dụng kết hợp giữa lưu lịch sử policy tại thời điểm nộp đơn và cơ chế tái tính điểm có kiểm soát, nhằm bảo đảm tính truy vết giữa quyết định phân bổ và căn cứ chính sách. Hướng tiếp cận này làm tăng khối lượng dữ liệu lưu trữ nhưng đổi lại giảm rủi ro tranh chấp, đồng thời tạo nền tảng cho các báo cáo kiểm định nội bộ.

\section*{3.2 Công nghệ backend}
Backend được triển khai trên nền tảng Node.js với framework Express.js, sử dụng MongoDB làm kho lưu trữ tài liệu và Mongoose làm ODM để biểu diễn schema và ràng buộc ở tầng ứng dụng. Node.js cung cấp môi trường thực thi JavaScript trên server phù hợp với hệ thống hiện có trong repository, cho phép tận dụng các tác vụ dựng chính sách, khởi tạo học kỳ và khôi phục xác thực hai lớp. Express.js đơn giản hóa việc tổ chức route và middleware cho các RESTful API cần thiết phục vụ các chức năng: allocation cycle, room assignment, policy inheritance và violation workflow. Các tuyến API tuân thủ phong cách REST, phân rạch rõ tài nguyên (users, rooms, applications, allocations) và trạng thái, thuận tiện cho việc kiểm thử tự động bằng Playwright và tích hợp với các tác vụ kiểm tra/tạo dữ liệu.

Ở mức thiết kế API, hệ thống tuân theo các nguyên tắc RESTful để đơn giản hóa hợp đồng giữa client và server: tài nguyên được xác định rõ ràng, sử dụng các phương thức HTTP phù hợp, và trả mã trạng thái hợp lệ cùng payload có cấu trúc chuẩn. Điều này giúp các script tự động và công cụ kiểm thử có thể tương tác trực tiếp bằng HTTP mà không cần phụ thuộc vào rendering layer. Bên cạnh đó, các lỗi nghiệp vụ được chuẩn hóa theo dạng JSON với mã lỗi có thể lược đồ hóa, hỗ trợ việc logging và phân loại lỗi khi chạy các allocation cycle lớn.

Về ghi log và giám sát, backend tích hợp logging có cấu trúc cho các hành động quan trọng (allocation, policy change, violation handling), kèm theo metadata để phục vụ truy vết. Thiết kế này giúp phân tích nguyên nhân khi phát sinh lỗi trong tác vụ phân phòng tự động hoặc khi migration làm lệch dữ liệu phân bổ.

Việc chọn MongoDB là phù hợp do mô hình dữ liệu có tính biến đổi và quan hệ một-nhiều không phức tạp: hồ sơ sinh viên, đơn đăng ký, và chính sách phân bổ có cấu trúc khác nhau giữa các học kỳ; mô hình document cho phép mở rộng trường mà không phá vỡ schema hiện hành. Mongoose hỗ trợ khai báo schema, các middleware kiểm tra tính hợp lệ khi save, và các phương thức truy vấn biểu diễn rõ ràng, từ đó tạo thuận lợi cho các cơ chế scoring và inheritance logic. Các thay thế cân nhắc gồm PostgreSQL hoặc hệ quản trị quan hệ khác; tuy nhiên lựa chọn đó sẽ làm tăng chi phí chuyển đổi mô hình dữ liệu và phức tạp khi xử lý tài liệu đa dạng (ví dụ các policy có trường động), đồng thời đòi hỏi phát triển ORM/DAO phức tạp hơn.

Về bảo mật và xác thực, backend sử dụng: lưu mật khẩu bằng bcrypt, xác thực bằng JWT cho API không trạng thái, và hỗ trợ TOTP/2FA cho các thao tác nhạy cảm như reset 2FA hoặc thay đổi thiết lập người dùng. Cơ chế này cân bằng giữa trải nghiệm người dùng và yêu cầu an toàn cho các thao tác quản trị. Nhóm tác vụ quản trị cho phép xử lý hàng loạt các trường hợp cần khôi phục 2FA theo quy trình kiểm soát quyền chặt chẽ.

\section*{3.3 Công nghệ frontend}
Frontend của hệ thống được triển khai theo kiến trúc đa trang (multi-page) sử dụng EJS làm template engine server-side, kết hợp HTML5, CSS3, JavaScript và Bootstrap để đảm bảo bố cục và tương tác người dùng cần thiết. Lựa chọn EJS xuất phát từ việc hệ thống cần các trang động phía server cho hai nhóm nghiệp vụ chính: khám phá phòng theo không gian và biểu mẫu đăng ký nâng cao, từ đó cho phép render dữ liệu phân bổ và trạng thái phòng ngay tại server, giảm phức tạp client và đơn giản hóa quy trình SEO/SSR cho trang quản trị.

Server-side rendering với EJS giúp duy trì tính nhất quán hiển thị cho các biểu mẫu nâng cao (enhanced application form), đồng thời giảm lượng logic cần chạy trên client khi thực hiện các thao tác liên quan đến policy inheritance hoặc tính toán điểm ưu tiên. Kiến trúc này phù hợp khi cần kiểm soát chặt chẽ luồng dữ liệu giữa server và client, tránh chi phí phát triển và vận hành của SPA, đồng thời tích hợp trực tiếp với middleware xác thực và các route RESTful. Một số công nghệ thay thế như React hoặc Vue (SPA) có lợi ở mặt trải nghiệm tương tác, nhưng sẽ yêu cầu thay đổi kiến trúc backend để phục vụ API nhiều hơn, tăng độ phức tạp đồng bộ trạng thái và không phù hợp với mục tiêu giữ đơn giản cho quy trình phân bổ và seeding hiện tại.

Trên lớp trình bày, hệ thống áp dụng nguyên tắc progressive enhancement: các trang cơ bản vẫn hoạt động khi JavaScript bị tắt, trong khi các tương tác nâng cao được bổ sung thông qua các module JavaScript nhỏ. Điều này cải thiện khả năng tương thích và giảm rủi ro lỗi trên các trình duyệt cũ, đồng thời giữ trải nghiệm mượt cho người dùng phổ thông. Hệ thống cũng chú ý đến khả năng truy cập (accessibility) và validation phía server để đảm bảo dữ liệu gửi lên luôn được kiểm tra trước khi xử lý.

\section*{3.4 Cơ sở dữ liệu}
Cơ sở dữ liệu chính của hệ thống là MongoDB; mô hình dữ liệu hướng document phù hợp với tính chất thay đổi của các thực thể: sinh viên, phòng, chính sách và đợt học (academic window). Một ví dụ mô tả cấu trúc tài liệu: collection \texttt{students} chứa các trường nhận dạng, \texttt{applications} lưu các đơn theo học kỳ kèm điểm ưu tiên (\texttt{priority\_score}), \texttt{rooms} lưu thông tin thuộc tính phòng và trạng thái, \texttt{policies} lưu cây kế thừa các quy tắc phân bổ. Việc biểu diễn priority score và lịch sử phân bổ dưới dạng tài liệu nhúng hoặc tham chiếu được cân nhắc tùy theo nhu cầu truy vấn: các truy vấn phân bổ chu kỳ thường cần truy cập nhanh tới dữ liệu đơn và policy, do đó một số trường tóm tắt được lưu trực tiếp trong document \texttt{applications} để giảm số lượng join.

Ở các truy vấn phân tích, hệ thống thường tính toán theo hai bước: thu thập dữ liệu thô và tạo các trường tóm tắt (precomputed fields) phục vụ thuật toán scoring, sau đó thực hiện sắp xếp và phân bổ. Cách tiếp cận này giúp giảm chi phí truy vấn phức tạp trong thời gian chạy allocation cycle và cho phép lưu lại lịch sử tính toán để phục vụ tra cứu và kiểm chứng sau khi phân bổ.

Ở tầng ứng dụng, Mongoose cung cấp schema validation và virtuals để xây dựng các thuộc tính tính toán (ví dụ \texttt{normalized\_score}), đồng thời hỗ trợ migration scripts khi cần tái cấu trúc dữ liệu. Các tác vụ seed dữ liệu phòng và dữ liệu hồ sơ mẫu thể hiện quy trình dựng dữ liệu test; tác vụ sinh dữ liệu mô phỏng quy mô lớn được dùng để thử nghiệm thuật toán phân bổ và kiểm thử tải.

Việc duy trì tập script seed và mock là một phần thiết yếu của quy trình phát triển: chúng cho phép tái tạo môi trường kiểm thử với các kịch bản biên (edge cases) như xung đột preference, thiếu thông tin sinh viên, hoặc thay đổi policy đột ngột. Điều này giúp giảm thiểu rủi ro khi triển khai thay đổi logic phân bổ lên môi trường thực tế.

Đối với bài toán policy inheritance, mô hình dữ liệu cần biểu diễn rõ quan hệ kế thừa theo cấp: quy định chung cấp trường, quy định theo khu nhà, và quy định theo loại phòng. Khi triển khai bằng MongoDB, cấu trúc document cho phép lưu cả tập luật kế thừa và tập luật ghi đè trong cùng một thực thể policy, từ đó giảm số vòng truy vấn khi cần tính kết quả cuối cùng cho từng hồ sơ. Tuy nhiên, nếu lạm dụng nhúng sâu, document sẽ khó bảo trì và khó kiểm tra nhất quán; do đó hệ thống chọn chiến lược kết hợp: phần cấu hình lõi được tham chiếu độc lập, phần snapshot phục vụ phân bổ được ghi lại theo học kỳ. Cách làm này cân bằng giữa hiệu quả truy vấn và khả năng kiểm soát thay đổi chính sách.

Một yêu cầu quan trọng khác là bảo đảm tính toàn vẹn khi chạy tác vụ xử lý vi phạm ở quy mô lớn. Các tác vụ loại này có thể cập nhật số lượng lớn hồ sơ trong thời gian ngắn; nếu không có cơ chế ghi nhận trạng thái trung gian, hệ thống dễ rơi vào tình huống dữ liệu hợp lệ một phần. Vì vậy, mỗi lần chạy tác vụ cần được ràng buộc bởi quy trình ba bước: sao lưu snapshot trước khi cập nhật, thực thi cập nhật theo lô với điều kiện rõ ràng, và đối soát hậu kiểm theo tập chỉ số (số hồ sơ chuyển trạng thái, số hồ sơ bị loại, số hồ sơ cần xem xét thủ công). Đây là điều kiện cần để dữ liệu phân bổ có thể dùng cho báo cáo quản trị.

\section*{3.5 Xác thực và bảo mật}
Chiến lược xác thực kết hợp sessionless token (JWT) cho API và cơ chế TOTP/2FA cho các thao tác nhạy cảm. Mật khẩu được băm bằng bcrypt với tham số salt phù hợp; các endpoint thay đổi dữ liệu nhạy cảm yêu cầu token có thời hạn ngắn và kiểm tra lại 2FA khi cần. Việc dùng JWT giúp đơn giản trong kịch bản phân tán và gọi API từ các script tự động, song cần lưu ý các rủi ro như token theft và replay — hệ thống bổ sung cơ chế refresh token và blacklist token khi rút quyền.
Thêm vào đó, về lưu trữ token ở phía client, hệ thống khuyến nghwwwwwwwwww
ị sử dụng cookie có thuộc tính HttpOnly và SameSite phù hợp cho access token/ngắn hạn, kết hợp refresh token được lưu an toàn trên server hoặc sử dụng cơ chế rotation để giảm nguy cơ bị đánh cắp. Hệ thống cũng triển khai kiểm tra hành vi bất thường (anomaly detection) trên các token đăng nhập để phát hiện hoạt động khả nghi.

Ở mức quyền truy cập, hệ thống áp dụng phân quyền theo vai trò (RBAC) cho các thao tác quản trị và kiểm soát truy cập theo tài nguyên (resource-level authorization) để đảm bảo nguyên tắc ít đặc quyền (least privilege). Mọi thao tác can thiệp dữ liệu nhạy cảm đều được audit và ghi log chi tiết.

Thêm vào đó, backend cần thực hiện các biện pháp phòng chống tấn công phổ biến: xác thực đầu vào (input validation), giới hạn tốc độ (rate limiting) trên các route quan trọng, và kiểm tra quyền truy cập ở nhiều mức (route-level và resource-level). Các tác vụ xử lý vi phạm sau điều tra cần vận hành dưới quyền hạn giới hạn cùng logging đầy đủ để phục vụ truy vết.

\section*{3.6 Công nghệ trực quan hóa và bản đồ}
Hệ thống sử dụng Leaflet.js để cung cấp bản đồ tương tác và hiển thị phân bố phòng theo không gian địa lý (ví dụ bản đồ khuôn viên hoặc sơ đồ tầng). Leaflet thích hợp cho các hiển thị lớp raster/vector nhẹ, cho phép tích hợp dễ dàng với dữ liệu phòng từ backend và tương tác để thực hiện live filtering hoặc chọn phòng. Việc sử dụng Leaflet được kết hợp với Bootstrap để duy trì trải nghiệm responsive trên nhiều thiết bị.

Trong thực tế tích hợp, dữ liệu phòng được cung cấp qua các dịch vụ REST dưới dạng GeoJSON hoặc các cấu trúc nhẹ khác; frontend tải các lớp (layers) theo yêu cầu và áp dụng clustering/tiling để tối ưu hiển thị khi có nhiều điểm. Các chiến lược tối ưu hóa gồm lazy loading cho các vùng bản đồ không hiển thị và caching các tile/vector đã tải để giảm băng thông.

Với module room explorer 360, tích hợp giữa tập ảnh/phông nền và siêu dữ liệu phòng cho phép xây dựng trải nghiệm tham quan. Khi mở rộng sang Three.js, cần xây dựng pipeline chuyển đổi mô hình từ Blender sang định dạng glTF phù hợp cho web, cùng các bước giảm độ phức tạp (mesh decimation) và nén tài sản để đảm bảo hiệu năng client.

Ngoài ra, hệ thống hiện có module room explorer dạng 360 để khám phá mô hình phòng; đây là phần trực quan hóa nâng cao cho phép người dùng duyệt nội thất phòng trước khi duyệt đơn. Là định hướng mở rộng, hệ thống đề xuất bổ sung module dựng hình bằng Blender theo quy trình: thu thập ảnh thực tế từng phòng, hiệu chỉnh tỉ lệ trong Blender, dựng mô hình nội thất chuẩn hóa, xuất tài sản 3D định dạng glTF và gắn mã định danh phòng tương ứng. Sau đó, module hiển thị 3D bằng Three.js nhận dữ liệu glTF để render tương tác trên trình duyệt. Cách làm này tăng khả năng trình bày không gian nhưng đồng thời phát sinh chi phí chuyển đổi dữ liệu và tối ưu hóa biểu diễn đối tượng. Ba lợi ích chính của hướng mở rộng này là: cải thiện tính trực quan cho người dùng, hỗ trợ kiểm tra vật lý không gian trước khi phân bổ, và tạo dữ liệu phong phú cho các công cụ phân tích; nhược điểm là yêu cầu bổ sung pipeline cho conversion, hosting tài sản 3D, và tối ưu hiệu suất phía trình duyệt.

\section*{3.7 Công nghệ kiểm thử và sinh dữ liệu}

Khung kiểm thử end-to-end sử dụng Playwright để mô phỏng hành vi người dùng trên giao diện multi-page và kiểm tra luồng phân bổ từ tạo đơn đến ghi nhận allocation. Playwright cho phép viết kiểm thử kịch bản tương tác phức tạp, bao gồm xác thực hai yếu tố và hành vi bất thường trong violation workflow. Song song, hệ thống duy trì nhóm tác vụ hỗ trợ sinh dữ liệu và seed cho chính sách, hồ sơ sinh viên và phân phòng tự động. Các tác vụ này đóng vai trò thiết yếu để tái tạo môi trường thử nghiệm cho các trường hợp cạnh (edge cases) của cơ chế scoring, cũng như để kiểm tra hoạt động sao lưu và khôi phục dữ liệu.

Chiến lược kiểm thử gồm nhiều lớp: unit test cho các hàm tính điểm và xử lý policy inheritance, integration test cho các service tương tác với MongoDB (sử dụng cơ sở dữ liệu tạm thời hoặc in-memory), và end-to-end test với Playwright trên workflow thực tế. Để đảm bảo độ che phủ và tính ổn định, pipeline CI thực hiện chạy đầy đủ các bước seed, tạo mock data bằng tác vụ sinh dữ liệu mô phỏng, rồi thực thi các kịch bản Playwright trên môi trường tĩnh hóa.

Ngoài ra, cần có bài kiểm thử tải và stress để đánh giá hiệu suất allocation cycle khi số lượng đơn lớn; các tập dữ liệu sinh tự động giúp mô phỏng các kịch bản này và xác định điểm nghẽn hệ thống trước khi triển khai trên môi trường sản xuất.

Ở góc độ tổ chức kiểm thử, hệ thống nên xác định rõ bộ chỉ số chất lượng cho từng nhóm chức năng thay vì chỉ dựa vào việc testcase chạy thành công. Với allocation cycle, chỉ số trọng tâm là độ ổn định của thứ tự ưu tiên khi dữ liệu đầu vào thay đổi nhỏ; với room assignment, chỉ số trọng tâm là tỉ lệ phân bổ thành công theo từng nhóm đối tượng và mức độ tuân thủ ràng buộc sức chứa. Đối với violation workflow, cần theo dõi thêm thời gian xử lý trung bình và tỉ lệ hồ sơ phải can thiệp thủ công sau khi script hoàn tất. Cách đo này giúp nhóm phát triển đánh giá chất lượng theo mục tiêu nghiệp vụ thay vì chỉ theo tiêu chí kỹ thuật thuần túy.

Một nội dung không thể bỏ qua là khả năng tái lập kết quả kiểm thử. Khi sử dụng tác vụ sinh dữ liệu mô phỏng, hệ thống nên quản lý seed sinh dữ liệu theo phiên bản để cùng một kịch bản có thể chạy lại và cho ra tập dữ liệu tương đương về phân bố thống kê. Nhờ đó, khi phát hiện sai lệch giữa hai lần chạy allocation, nhóm phát triển có thể khoanh vùng nguyên nhân nhanh hơn: do thay đổi thuật toán, do thay đổi policy, hay do thay đổi dữ liệu. Kết hợp với nhật ký thực thi tác vụ và báo cáo kiểm thử tự động, quy trình này tạo nền tảng vững chắc cho bảo trì dài hạn và nâng cấp chức năng trong các học kỳ tiếp theo.

Triển khai pipeline kiểm thử nên bao gồm: (1) chạy các migration và seed cơ sở dữ liệu, (2) thực thi các script tạo mock data, (3) chạy test suite Playwright trên các kịch bản phân bổ chính và xử lý vi phạm, (4) xác minh tính toàn vẹn dữ liệu sau khi chạy các script tự động. Việc tự động hóa chuỗi này giúp phát hiện sớm lỗi logic trong policy inheritance hoặc sai lệch trong scoring mechanism.

\bigskip
Các nguồn tài liệu chính được sử dụng trực tiếp trong chương này gồm: Node.js Documentation, Express Documentation, MongoDB Manual, Mongoose Documentation, Bootstrap Documentation, Leaflet Documentation, Playwright Documentation và Three.js Documentation.
